\section{The first CUE heat collision is an extreme-gap event}
\label{sec:analytic-cue-first-collision}

Let $0\le\theta_1<\cdots<\theta_N<2\pi$ be the eigenangles
of Haar $U(N)$, and extend them periodically by
$\theta_{j+N}=\theta_j+2\pi$. The real characteristic
trigonometric polynomial is
\[
 Z_N(z)=\prod_{j=1}^N2\sin\frac{z-\theta_j}{2},\qquad
 Z_{N,-a}=e^{a\partial_z^2}Z_N .
\]
For odd $N$ it is antiperiodic, and its zero divisor is
still $2\pi$-periodic. Negative heat time has exactly the
same sign convention as $H_{-a}$ in the arithmetic family.
Define
\[
 \delta_N=\min_j(\theta_{j+1}-\theta_j),\qquad
 A_N=\inf\{a>0:Z_{N,-a}\text{ has a nonreal zero}\}.
\]
This section concerns this specified CUE model, with no
arithmetic transfer assumption.

\begin{theorem}[NN-A024: first CUE collision law]
For $N\ge2$, $A_N$ is positive and finite almost surely.
As $N\to\infty$,
\[
                    \frac{8A_N}{\delta_N^2}
                                  \longrightarrow1
                                  \quad\hbox{in probability}.
\]
More precisely the ratio is $1+O_{\mathbb P}(N^{-1/3})$.
With probability tending to one, the first collision is
the one continued from the initially smallest cyclic gap.
For $u\ge0$,
\begin{equation}\label{eq:analytic-cue-first-time-law}
 \lim_{N\to\infty}
 \mathbb P\left\{
       \frac{8N^{8/3}A_N}{(72\pi)^{2/3}}>u\right\}
                              =e^{-u^{3/2}}.
\end{equation}
Equivalently, after unfolding the mean spacing to one,
the first negative heat time is of order $N^{-2/3}$,
smaller than the mean-spacing-squared scale, which is one.
\end{theorem}

\begin{proof}
We first prove a deterministic global lower bound for
the first loss of reality. Write
\[
 Z_N(z)=\sum_{k=0}^Nc_ke^{i(k-N/2)z}.
\]
The coefficients $c_0,c_N$ are nonzero. Under
$e^{a\partial_z^2}$ the $k$th coefficient is multiplied
by $e^{-a(k-N/2)^2}$. After multiplication by
$e^{iNz/2}$ this is a degree-$N$ polynomial in $w=e^{iz}$.
Its leading and constant coefficients remain nonzero,
and on a compact $a$ interval the normalized polynomial
coefficients and their reciprocals at the two ends stay
bounded. The $N$ roots therefore remain in a compact
annulus in $\mathbb C^*$.

While all these roots are on the unit circle and simple,
their angles have real analytic continuations $y_j(a)$.
There is no loss of a zero at imaginary infinity at a
finite time; and a simple real zero remains real by the
implicit function theorem and conjugation. Thus a first
loss of reality must occur at a finite real collision.
Before that collision the factorization into real sine
factors gives
\begin{equation}\label{eq:analytic-circular-backward-ode}
 y_i'(a)=-\sum_{j\ne i}\cot\frac{y_i(a)-y_j(a)}2.
\end{equation}
For an empty adjacent arc of length
$g_i=y_{i+1}-y_i\in(0,2\pi)$, the two mutual terms give
$-2\cot(g_i/2)$ to $g_i'$. Every remaining cotangent
difference is negative before the minus sign: along
this empty arc the cotangent is strictly decreasing
and crosses no pole. Therefore
\[
                  g_i'(a)\ge-2\cot(g_i(a)/2),
                  \qquad (g_i^2)'(a)\ge-8.
\]
There is a slightly sharper integrated form:
\[
       \frac d{da}\cos(g_i(a)/2)
                       \le\cos(g_i(a)/2).
\]
If this gap closes at $A$, its left side integrates to
$e^{-A}\le\cos(g_i(0)/2)$. Since
$g_i(0)\ge\delta_N$ and $\delta_N<\pi$ almost surely,
\begin{equation}\label{eq:analytic-first-time-lower}
       A_N\ge-\log\cos(\delta_N/2)\ge\delta_N^2/8.
\end{equation}
This rules out an earlier collision anywhere on the
circle, not merely near the smallest pair.

For completeness, the first collision is almost surely
finite. When $N$ is even the constant Fourier coefficient
is almost surely nonzero and eventually dominates all
others, so for large $a$ there are no real zeros.
When $N$ is odd, the almost surely nonzero frequencies
$\pm1/2$ eventually dominate, leaving one simple real
zero modulo $2\pi$, fewer than $N$ for $N\ge3$.
Uniform $C^1$ convergence after rescaling gives this
last zero count. Each asserted Fourier coefficient is
a nonzero real analytic function of the eigenangles:
one can check nonvanishing when the angles are all
close to zero, where the limiting function is
$(2\sin(z/2))^N$. Its vanishing set has Lebesgue
measure zero, and the CUE law has an absolutely
continuous density. The preceding annulus argument
then supplies a finite first collision.

We next prove a uniform isolation estimate for the
smallest pair. Put $\ell_N=N^{-7/6}$. With probability
tending to one, every circular interval of length
$\ell_N$ contains at most two eigenangles.
Here is a direct determinantal proof. Let
\[
 v_N(x)=(1,e^{ix},\ldots,e^{i(N-1)x})/\sqrt{2\pi}.
\]
The three-point intensity is
$\rho_3(x_1,x_2,x_3)=
 \|v_N(x_1)\wedge v_N(x_2)\wedge v_N(x_3)\|^2$.
Successive divided differences give
\[
 \rho_3(x_1,x_2,x_3)
       \le C N^9\prod_{i<j}|x_i-x_j|^2.
\]
Indeed the norms of $v_N,v_N',v_N''$ are respectively
$O(N^{1/2}),O(N^{3/2}),O(N^{5/2})$. The determinant
of the three divided-difference columns extracts the
Vandermonde product, and their norms are bounded by
these derivative norms (with a factor $1/2$ in the
second divided difference). This applies to real
lifts of any short arc, uniformly around the circle.
The expected number of triples lying in a forward
arc of length $\ell_N$ is therefore at most
\[
 C N^9\ell_N^8=O(N^{-1/3}).
\]
If such an arc contains three points, one of its points
can be chosen as the start of a forward arc containing
the other two. Markov's inequality proves the claim.

The classical extreme-gap theorem of Ben Arous--Bourgade
gives
\begin{equation}\label{eq:analytic-CUE-min-gap}
 \frac{N^{4/3}\delta_N}{(72\pi)^{1/3}}
          \ \Longrightarrow\ Y,\qquad
                    \mathbb P\{Y>y\}=e^{-y^3}\quad(y\ge0).
\end{equation}
In particular $\delta_N=O_{\mathbb P}(N^{-4/3})$ and
$\delta_N/\ell_N\to0$ in probability.
Choose the minimum pair, using real lifts
$b-\delta_N/2,b+\delta_N/2$. On the preceding isolation
event, for $\delta_N\le\ell_N/2$ every other zero on the
periodic real line has distance at least $\ell_N/2$
from $b$. Each interval of length $\ell_N/2$ contains
at most two zeros. Summing successive distance shells
therefore gives
\[
 S_N=\sum_{\substack{Z_N(\xi)=0\\
                   \xi\ne b-\delta_N/2,b+\delta_N/2}}
                     (b-\xi)^{-2}
                          \le C\ell_N^{-2}.
\]
The periodic copies are included in this sum. Thus the
pair parameter satisfies
\[
                    \kappa_N=\delta_N^2S_N
                             =O_{\mathbb P}(N^{-1/3}).
\]
This includes the exceptional probability
$O(N^{-1/3})$ of the isolation event failing:
outside that event no bound is asserted, while its
probability tends to zero, as required for the stated
bound in probability.

The real entire function $Z_N$ has order one and only
real zeros. Removing the selected pair leaves an
order-one function with the required canonical product.
NN-A023 applies, including its unrestricted real drift.
It gives an actual collision of that pair at
\[
 a_{\min}=\frac{\delta_N^2}{8}
                    \{1+\kappa_N/4+O(\kappa_N^2)\}.
\]
The first global collision can occur no later than this
one. Combine this upper bound with
\eqref{eq:analytic-first-time-lower} to obtain
$8A_N/\delta_N^2=1+O_{\mathbb P}(N^{-1/3})$.

To identify the colliding pair, let $\delta_N^{(2)}$
be the second-smallest initial gap. The same source's
joint extreme-gap theorem implies that
$\delta_N^{(2)}/\delta_N$ converges in distribution
to a variable strictly larger than one almost surely.
Since $\kappa_N\to0$ in probability, with probability
tending to one
\[
                    (\delta_N^{(2)})^2/8>a_{\min}.
\]
Every other initially labelled cyclic gap has its
first possible closing time at least its initial
squared length divided by eight, by the universal
gap inequality. Hence none can collide first before
$a_{\min}$, proving the assertion about the smallest pair.
Finally combine \eqref{eq:analytic-CUE-min-gap} with
$8A_N/\delta_N^2\to1$ and the continuous mapping theorem:
the variable in \eqref{eq:analytic-cue-first-time-law}
converges to $Y^2$, whose survival function is
$e^{-u^{3/2}}$.
\end{proof}

\paragraph{Exact two-angle normalization check.}
For just two angles $b\pm\delta/2$,
\[
 Z_2(z)=2\cos(\delta/2)-2\cos(z-b),\qquad
 Z_{2,-a}(z)=2\cos(\delta/2)-2e^{-a}\cos(z-b).
\]
Consequently
\[
              A_2=-\log\cos(\delta/2)
                  =\delta^2/8+\delta^4/192+O(\delta^6)
\]
when $0<\delta<\pi$. The complementary periodic zero
sum in NN-A023 is
\[
          S=\frac1{2\sin^2(\delta/4)}-\frac8{\delta^2}
                                  =\frac16+O(\delta^2).
\]
Its correction $\delta^2\kappa/32$ is exactly
$\delta^4/192$ to first order. This checks both the
heat-sign normalization and the interaction coefficient
by a closed-form example.

\paragraph{Primary locator and arithmetic scope.}
The only imported random-matrix asymptotic is
Ben Arous--Bourgade, \emph{Extreme gaps between eigenvalues
of random matrices}, Annals of Probability~41 (2013),
2648--2681, DOI~10.1214/11-AOP710, Theorem~1.1
printed p.~2650 and Corollary~1.2 p.~2651.
Their full-circle gap intensity is
$u^2\,du/(24\pi)$ for the scale $N^{4/3}\delta$,
giving the constant $72\pi$ in
\eqref{eq:analytic-CUE-min-gap}. The no-triples estimate,
heat dynamics, global earliest-collision bound and
passage from gaps to collision times were proved above.
The theorem uses neither A nor M.

The model also distinguishes two different heuristic
questions in the supplied consultation. Matching
$\sqrt a$ to the mean arithmetic gap gives a
\emph{typical-gap} or bulk scale
$\log X\asymp a^{-1/2}$. It is not the scale of the
first rare collision in the CUE model. If one
additionally postulates the analogous cubic small-gap
Poisson law for arithmetic zeros, together with
isolation and exhaustion of early collisions by
those pairs, the heat-coordinate density
$D(x)\sim\log x/(4\pi)$ predicts a count
\[
 \frac{\pi^2}{9}(8a)^{3/2}\int_0^X D(x)^4\,dx
           \sim\frac{\sqrt2}{144\pi^2}
                              a^{3/2}X(\log X)^4.
\]
The first-event scale in that stronger model is
$X\asymp a^{-3/2}(\log(1/a))^{-4}$, rather than an
exponential scale. This displayed arithmetic conversion
is explicitly a conditional model prediction.
Neither A, Montgomery's restricted pair-correlation
theorem, nor a fixed-scale sine-kernel limit supplies
the required extreme-gap process or the exhaustion
of all early arithmetic collisions. NN-A023 proves
the local dynamics once the isolation parameter is small.

\paragraph{An actual algebraic family and its discriminant.}
Put $q=e^{-a/4}$ and $w=e^{iz}$. The preceding Fourier
formula gives the polynomial
\[
              P_N(q,w)=
                    \sum_{k=0}^Nc_kq^{(2k-N)^2}w^k
                 \in\mathbb C[q,q^{-1},w,w^{-1}].
\]
Its leading and constant coefficients are units over
$\mathbb C[q,q^{-1}]$. Dividing by the leading
coefficient exhibits
\[
 \operatorname{Spec}
       \mathbb C[q,q^{-1},w,w^{-1}]/(P_N)
                      \longrightarrow\mathbb G_{m,q}
\]
as a finite free family of degree $N$; inversion of $w$
is automatic because its constant coefficient is a unit.
It is \'{e}tale precisely off its polynomial
discriminant, by the derivative or trace-pairing
criterion. The exact Grothendieck locator is
SGA~1, Expos\'e~I, section~4, Proposition~10
(invertibility of the discriminant for a finite
locally free cover), with Definition~1 in that
section and section~3, Proposition~1 for the
flat/unramified and differential formulations.
At the first collision identified above with
probability tending to one, the local rank-two
component has a simple discriminant zero and
ramified double-cover normalization as in NN-A023.
The circle is the fixed locus of
$w\mapsto1/\overline w$ on the real parameter ray.
This is a genuine complex algebraic zero family
for each CUE sample; it is not an arithmetic
Frobenius construction or a canonical-basis model.
